\chapter{Introducción}

\section{El cáncer como sistema físico fuera del equilibrio}
La progresión tumoral puede interpretarse como un sistema complejo, abierto y fuera del equilibrio, con flujos de materia y energía que mantienen estados metaestables. Esta visión permite reutilizar herramientas de física estadística y de sistemas dinámicos para describir transiciones entre coexistencia y dominancia tumoral.

\section{Limitaciones de modelos clásicos}
Modelos logísticos o Lotka--Volterra simétricos no capturan umbrales críticos ni asimetrías en las interacciones cáncer--tejido sano--sistema inmune. Tampoco incorporan respuestas no lineales del microambiente ni la posibilidad de histeresis o patrones espaciales.

\section{Efecto Allee y no reciprocidad}
El efecto Allee introduce umbrales de viabilidad en poblaciones bajas; la no reciprocidad recoge que las acciones de una población sobre otra no son necesariamente simétricas. Su combinación produce bifurcaciones ricas y múltiples ramas de equilibrio, clave para entender latencia tumoral y despegues abruptos.

\section{Control inmunológico como campo externo}
La inmunoterapia se modela aquí como un campo de control $u(x,y,t)$ o como un parámetro binario $\mu$ que activa términos variacionales. Este enfoque permite comparar dinámicas puramente biológicas con dinámicas guiadas por energía libre tipo modelo C.

\section{Objetivos y contribuciones}
\begin{itemize}
    \item Formular un modelo ODE/PDE con efecto Allee (débil/fuerte), interacciones no recíprocas y control inmunológico.
    \item Caracterizar estados estacionarios, su estabilidad y bifurcaciones relevantes.
    \item Analizar dinámica espacio--temporal y patrones usando simulaciones numéricas.
    \item Evaluar estrategias de control inmunológico y su impacto en la extinción o contención del tumor.
\end{itemize}

\noindent\textbf{Contribuciones a la frontera (física y matemática).} (i) Integración de un marco variacional tipo modelo C (Halperin--Hohenberg) con dinámicas cáncer--tejido sano--inmune y parámetro $\mu$ que modifica el funcional efectivo. (ii) Cuantificación de organización mesoscópica mediante longitudes de correlación $\xi(t)$ y coarsening subdifusivo $\xi(t)\sim e^{\alpha}t^{1/2}$ en el horizonte simulado. (iii) Separación conceptual entre el papel de $\mu$ como \emph{selector morfológico} (espectro espacial, fragmentación) y el de $u(x,y,t)$ como \emph{intervención geométrica} sobre coherencia tumoral y acoplamiento cáncer--inmune. El desarrollo cuantitativo aparece en los caps.~4--6 y la interpretación física en el Cap.~\ref{chap:discusion}.

\noindent\textbf{Contribuciones interpretativas (biología y control).} (i) Lectura de umbrales de Allee como separatrices inestables que organizan rutas transitorias, en lugar de atractores estables, en los barridos homogéneos reportados. (ii) Uso de métricas $c_c$ y $c_i$ como lenguaje reproducible para comparar protocolos de inmunoterapia modelada. (iii) Contraste sistemático Allee débil vs fuerte y su ``dureza'' geométrica respecto al tamaño de dominios. Estas lecturas son \emph{hipótesis de interpretación} del modelo; las limitaciones y el grado de evidencia de cada afirmación se discuten en el Cap.~\ref{chap:discusion}.



















